
\documentclass{article}
\usepackage{colortbl}
\usepackage{makecell}
\usepackage{multirow}
\usepackage{supertabular}

\begin{document}

\newcounter{utterance}

\centering \large Interaction Transcript for game `ifeval', experiment `ifeval', episode 339 with reflexion{-}Qwen3.5{-}27B.
\vspace{24pt}

{ \footnotesize  \setcounter{utterance}{1}
\setlength{\tabcolsep}{0pt}
\begin{supertabular}{c@{$\;$}|p{.15\linewidth}@{}p{.15\linewidth}p{.15\linewidth}p{.15\linewidth}p{.15\linewidth}p{.15\linewidth}}
    \# & \multicolumn{2}{c}{Player} && \multicolumn{2}{c}{Game Master} \\
    \hline

    \theutterance \stepcounter{utterance}  
    & & & \multicolumn{4}{p{0.6\linewidth}}{
        \cellcolor[rgb]{0.9,0.9,0.9}{
            \makecell[{{p{\linewidth}}}]{
                \texttt{\tiny{[P1$\langle$GM]}}
                \texttt{You have attempted this same task before and did not succeed. These are the plans you wrote after those attempts.} \\
\texttt{Trial \#1: New plan: I failed to meet the specific formatting requirement of highlighting "at least 3 sections" using the requested markdown style. In my previous attempt, I used standard Markdown headers (`\#\#\#`) and bold text (`**`), but I did not use the specific italicized format (`*text*`) to highlight distinct sections as explicitly demonstrated in the prompt's example (`*highlighted section*`). Additionally, while I included a placeholder, the primary failure was the highlighting strategy. Next time, I will ensure that I wrap the titles or key introductory phrases of at least three distinct sections in single asterisks (e.g., `*Section Title*`) to strictly adhere to the "highlighted section" format requested, rather than relying on headers or bolding. I will also verify that the placeholder `{[}address{]}` remains present.} \\
\\ 
\texttt{Trial \#2: New plan: I failed to strictly follow the formatting instruction to "highlight at least 3 sections" using the specific example style `*highlighted section*`. In my previous attempt, I used standard Markdown headers (`\#\#\#`) and bold text (`**`), which did not match the requested single{-}asterisk italicized format for the section highlights. Although I included a placeholder, the formatting of the sections was the primary cause of failure. Next time, I will explicitly wrap the titles or key phrases of at least three distinct sections in single asterisks (e.g., `*The Outer Courtyard*`, `*The Main Hall*`, `*The Worship Hall*`) to match the prompt's example exactly, while ensuring the placeholder `{[}address{]}` is still included in the text.} \\
\\ 
\\ 
\texttt{What is inside Shinto shrines? Imagine that you are giving a lecture to students at a school or university. Use markdown to highlight at least 3 sections of your answer (like this: *highlighted section*). Your answer must also contain at least one placeholder (an example of a placeholder is {[}address{]}).} \\
            }
        }
    }
     \\ \\

    \theutterance \stepcounter{utterance}  
    & \multicolumn{4}{p{0.6\linewidth}}{
        \cellcolor[rgb]{0.561,0.737,0.561}{
            \makecell[{{p{\linewidth}}}]{
                \texttt{\tiny{[P1$\rangle$GM]}}
                \texttt{Welcome, everyone, to today's lecture on the architecture and spiritual layout of Shinto shrines. As we explore these sacred spaces, it is important to understand that unlike many other religious structures, Shinto shrines are not primarily designed for congregational worship inside a large building. Instead, they are carefully curated environments meant to facilitate a connection between the human world and the *kami* (spirits or deities).} \\
\\ 
\texttt{*The Outer Courtyard and Purification*} \\
\texttt{Upon entering a shrine, you will first encounter the *sando*, or approach path, which often leads to a *torii* gate marking the transition from the mundane to the sacred. Before approaching the main structures, visitors typically stop at the *temizuya*, a water pavilion. Here, the act of *temizu* is performed, where one ritually washes their hands and mouth to purify themselves physically and spiritually. This area is crucial as it sets the tone for the visit, emphasizing cleanliness and respect before entering the inner sanctum.} \\
\\ 
\texttt{*The Main Hall and the Sacred Presence*} \\
\texttt{Moving past the purification area, you arrive at the *honden*, or main hall. This is the most sacred part of the shrine, yet it is often the most inaccessible to the general public. Inside the *honden*, the *shintai* (the physical object that houses the *kami*) is enshrined. This object can be a mirror, a sword, a jewel, or a natural stone, but it is rarely displayed openly. The *honden* is usually a simple, unadorned structure, reflecting the Shinto emphasis on purity and the natural world rather than elaborate iconography.} \\
\\ 
\texttt{*The Worship Hall and Community Interaction*} \\
\texttt{While the *honden* is reserved for the *kami*, the *haiden* is the hall where humans gather to pray. Located in front of the *honden*, this is the space where visitors offer prayers, shake the *suzu* (bell) to call the attention of the *kami*, and make offerings. It is here that the community aspect of Shinto is most visible, as people from all walks of life come together to give thanks or ask for blessings. If you are planning a visit to a specific shrine, you might want to check the schedule at {[}address{]} to see if there are any special festivals or ceremonies taking place.} \\
\\ 
\texttt{In summary, the layout of a Shinto shrine guides the visitor through a journey of purification, reverence, and prayer, with each section serving a distinct spiritual function. Thank you for your attention, and I look forward to our discussion.} \\
            }
        }
    }
    & & \\ \\

    \theutterance \stepcounter{utterance}  
    & & & \multicolumn{2}{p{0.3\linewidth}}{
        \cellcolor[rgb]{0.9,0.9,0.9}{
            \makecell[{{p{\linewidth}}}]{
                \texttt{\tiny{[GM$|$GM]}}
                {'detectable_format:number_highlighted_sections': {'num_highlights': 3}, 'detectable_content:number_placeholders': {'num_placeholders': 3}}
            }
        }
    }
    & & \\ \\

    \theutterance \stepcounter{utterance}  
    & & & \multicolumn{2}{p{0.3\linewidth}}{
        \cellcolor[rgb]{0.9,0.9,0.9}{
            \makecell[{{p{\linewidth}}}]{
                \texttt{\tiny{[GM$|$GM]}}
                \texttt{game\_result = LOSE} \\
            }
        }
    }
    & & \\ \\

    \theutterance \stepcounter{utterance}  
    & & \multicolumn{4}{p{0.6\linewidth}}{
        \cellcolor[rgb]{0.95,0.95,0.95}{
            \makecell[{{p{\linewidth}}}]{
                \texttt{\tiny{[P1$\leftrightarrow$P]}}
                \texttt{Post{-}episode Reflexion (not sent to GM):} \\
\texttt{New plan: I failed to interpret the instruction "highlight at least 3 sections" correctly. The prompt provided a specific example format `*highlighted section*`, which implies wrapping the entire content of a section (or a significant block of text) in single asterisks to create an italicized block, rather than just italicizing a title or using standard headers. In my previous attempt, I used headers and only italicized specific terms within the text, which did not constitute "highlighting a section" as demonstrated. Next time, I will structure the response so that three distinct paragraphs or blocks of text are fully enclosed in single asterisks (e.g., `*This entire paragraph is a highlighted section...*`) to strictly match the visual style of the example. I will also ensure the placeholder `{[}address{]}` is included within one of these highlighted sections or the surrounding text.} \\
            }
        }
    }
    & \\ \\

\end{supertabular}
}

\end{document}
